\chapter{Wymierna kontrakcja: od grafów plumbingowych do egzotycznych czterorozmaitości}
\label{ch:tg-rational-blowdown}

Przy zwykłym \emph{rozdmuchaniu} (\emph{blow-up}) gładkiej
czterorozmaitości wprowadzamy sferę o samoprzecięciu $-1$,
a przy kontrakcji (\emph{blow-down}) taką sferę usuwamy.
Istnieje istotnie subtelniejsza operacja: można usunąć całe
\emph{ujemnie określone plumbingowanie sfer} i zastąpić je
czterorozmaitością, która ma wymierną homologię punktu.
Nazywa się to \emph{wymierną kontrakcją} (\emph{rational
blow-down}). Wprowadzona przez Fintushela i Sterna operacja
dostarcza metod konstruowania egzotycznych struktur gładkich.

Kluczowe jest rozróżnienie trzech kroków: (i) obliczamy brzeg
plumbingowania, (ii) korzystamy z istnienia \emph{innego}
wypełnienia tego samego brzegu, (iii) dopiero wtedy
przeprowadzamy operację wycięcia i sklejenia. Rachunek
macierzy przecięcia uzasadnia pierwszy krok homologicznie,
ale sam nie dowodzi istnienia drugiego wypełnienia.

\section{Łańcuch \texorpdfstring{$C_p$}{Cp} i jego ujemnie określona forma przecięcia}
\label{sec:tg-rb-config}

\begin{definicja}[Konfiguracja Fintushela--Sterna]
Niech $p\geq2$ będzie liczbą całkowitą.
Przez $C_p$ oznaczamy zwartą, spójną, zorientowaną,
gładką czterorozmaitość plumbingową związaną z liniowym
grafem o $p-1$ wierzchołkach, z wagami
$-(p+2),-2,\ldots,-2$; waga $-2$ występuje
$p-2$ razy, a kolejne krawędzie mają dodatni znak
przecięcia. W bazie klas sfer będących przekrojami zerowymi,
zgodnie z opisem formy przecięcia plumbingu drzewiastego,
forma przecięcia jest opisana (twierdzenie
\ref{thm:tg-plumbing-boundary} i opis formy w rozdziale
o grafach plumbingowych)
macierzą symetryczną $Q_p$ stopnia $p-1$, której elementy
określamy wzorem
\[
(Q_p)_{ij}=\begin{cases}
-(p+2),&i=j=1,\\
-2,&i=j\geq2,\\
1,&|i-j|=1,\\
0,&\text{w pozostałych przypadkach}.
\end{cases}
\]
W szczególności $Q_2=[-4]$,
$Q_3=\begin{psmallmatrix}-5&1\\1&-2\end{psmallmatrix}$.
Dla $p\geq4$ jest to macierz trójdiagonalna,
z początkowym wpisem $-(p+2)$, dalszymi wpisami
$-2$ i jedynkami bezpośrednio przy przekątnej.
\end{definicja}

\begin{center}
\begin{tikzpicture}[>=Stealth,scale=.95]
  \node[tgRedVertex,minimum size=10mm] (a) at (-3.35,0) {$-(p+2)$};
  \node[tgBlueVertex] (b) at (-1.35,0) {$-2$};
  \node[tgBlueVertex] (c) at (.20,0) {$-2$};
  \node[font=\large,text=tgNavy] (d) at (1.43,0) {$\cdots$};
  \node[tgBlueVertex] (e) at (2.80,0) {$-2$};
  \draw[tgTree] (a)--(b)--(c)--(d)--(e);
  \node[font=\footnotesize,text=tgNavy] at (0,-.85)
    {$p-1$ wierzchołków; rysunek schematyczny dla $p\geq6$.};
\end{tikzpicture}
\end{center}

\begin{lemat}[Minory główne i określoność]
\label{lem:tg-rb-minors}
Dla każdej liczby całkowitej $p\geq2$ i każdego
$1\leq k\leq p-1$ wyznacznik górnego lewego bloku
$Q_p$ rozmiaru $k$ jest równy
$(-1)^k(1+k(p+1))$. W szczególności
\[
\det Q_p=(-1)^{p-1}p^2,
\qquad \operatorname{sign}(Q_p)=-(p-1).
\]
\end{lemat}

\begin{proof}
Niech $d_k$ będzie wyznacznikiem $k\times k$
początkowego bloku macierzy $-Q_p$.
Nie zakładamy z góry dodatniości $d_k$.
Przyjmujemy $d_0=1$. Mamy $d_1=p+2$.
Dla $2\leq k\leq p-1$ rozwinięcie wyznacznika
trójdiagonalnego względem ostatniego wiersza daje
$d_k=2d_{k-1}-d_{k-2}$:
składnik diagonalny wnosi $2d_{k-1}$,
a składnik odpowiadający parze wpisów $-1$
na sąsiednich pozycjach daje $-d_{k-2}$.
Dowodzimy indukcyjnie,
że $d_k=1+k(p+1)$: sprawdza się to dla $k=0,1$,
a jeśli wzór zachodzi dla dwóch poprzednich indeksów,
to $2[1+(k-1)(p+1)]-[1+(k-2)(p+1)]
=1+k(p+1)$.
Każdy początkowy minor główny macierzy $-Q_p$
jest dodatni, więc kryterium Sylvestera dowodzi,
że $-Q_p$ jest dodatnio określona. Wstawienie
$k=p-1$ daje $d_{p-1}=1+(p-1)(p+1)=p^2$.
Dla $p=2$ rozumowanie obejmuje pojedynczy minor $4$.
\end{proof}

\begin{twierdzenie}[Homologia brzegu $C_p$]
\label{thm:tg-rb-h1boundary}
Dla każdego $p\geq2$ brzeg $Y_p:=\partial C_p$
jest wymierną sferą homologiczną i zachodzi
$H_1(Y_p;\mathbb Z)\cong\mathbb Z/p^2\mathbb Z$.
\end{twierdzenie}

\begin{proof}
Twierdzenie o brzegu drzewa plumbingowego daje
$H_1(Y_p;\mathbb Z)\cong\operatorname{coker}Q_p$.
Z lematu \ref{lem:tg-rb-minors} znamy już rząd
tej grupy, równy $p^2$. Pokażemy, że jest cykliczna;
nie wynika to z samego wyznacznika.

Niech $e_1,\ldots,e_{p-1}$ oznaczają generatory
prezentacji ilorazowej, w kolejności od wierzchołka
o wadze $-(p+2)$. Jeśli $p\geq3$, pierwsza relacja
ma postać $-(p+2)e_1+e_2=0$, a kolejne relacje
wewnętrzne dają $e_{j+1}=2e_j-e_{j-1}$.
Stąd przez indukcję
$e_j=[1+(j-1)(p+1)]e_1$ dla $1\leq j\leq p-1$.
Ostatnia relacja $e_{p-2}-2e_{p-1}=0$
(dla $p=3$ jest to druga z dwóch relacji;
wtedy nie występują relacje wewnętrzne) staje się
$p^2e_1=0$, z pominięciem nieistotnego znaku.
Istotnie, współczynnik ostatniej relacji jest
$1+(p-3)(p+1)-2[1+(p-2)(p+1)]=-p^2$.
Tak więc cała grupa jest generowana przez $e_1$,
a jej rząd wynosi $p^2$. Dla $p=2$ jedyna relacja
$-4e_1=0$ daje ten sam wynik.
Ponieważ $Y_p$ jest zamkniętą, spójną,
zorientowaną trójrozmaitością, ze skończoności
$H_1(Y_p;\mathbb Z)$ wynika $H_1(Y_p;\mathbb Q)=0$.
Dualność Poincarégo nad $\mathbb Q$ i uniwersalne
twierdzenie o współczynnikach dają
$H_2(Y_p;\mathbb Q)=0$, natomiast
$H_0(Y_p;\mathbb Q)=H_3(Y_p;\mathbb Q)=\mathbb Q$.
To uzasadnia również część tezy dotyczącą
wymiernej sfery homologicznej.
\end{proof}

\section{Brzeg jako przestrzeń soczewkowa}
\label{sec:tg-rb-lens}

\begin{lemat}[Ułamek łańcuchowy konfiguracji]
\label{lem:tg-rb-cfrac}
Dla każdej liczby całkowitej $p\geq2$ zachodzi
\[
[p+2,\underbrace{2,\ldots,2}_{p-2}]^-
=\frac{p^2}{p-1}.
\]
\end{lemat}

\begin{proof}
Przez $[a_1,\ldots,a_k]^-$ oznaczamy
$a_1-1/(a_2-1/(\cdots-1/a_k))$.
Z indukcji otrzymujemy $[\underbrace{2,\ldots,2}_{m}]^-
=(m+1)/m$ dla $m\geq1$: dla $m=1$ jest to $2$,
a dodanie kolejnego wyrazu $2$ z lewej daje
$2-m/(m+1)=(m+2)/(m+1)$.
Gdy $p\geq3$, ogon ułamka ma $p-2$ dwójki,
a więc wartość $(p-1)/(p-2)$. Otrzymujemy
\[
[p+2,2,\ldots,2]^-
=p+2-\frac{p-2}{p-1}
=\frac{(p+2)(p-1)-(p-2)}{p-1}
=\frac{p^2}{p-1}.
\]
Dla $p=2$ lewa strona to $[4]^-=4$.
\end{proof}

\begin{twierdzenie}[Identyfikacja brzegu]
\label{thm:tg-rb-lens}
Przy konwencji, w której $L(n,q)$ oznacza wynik
$n/q$-chirurgii na dodatnio zorientowanym węźle
trywialnym w $S^3$, brzeg zorientowanej
czterorozmaitości $C_p$ jest orientacyjnie
dyfeomorficzny z $L(p^2,1-p)$.
\end{twierdzenie}

\begin{proof}
Dla $p=2$ diagram składa się z jednego
węzła trywialnego o ramowaniu $-4$ i teza
wynika z definicji $L(4,-1)$.
Dla $p\ge3$ stosujemy geometryczną redukcję
ramion plumbingu gwiaździstego z twierdzenia
\ref{thm:tg-star-seifert}.
Za centrum przyjmujemy wierzchołek o wadze
$b=-(p+2)$, a pozostały łańcuch jest jednym
ramieniem o ułamku $\alpha/\beta=(p-1)/(p-2)$.
Po redukcji otrzymujemy iloczyn pierścienia
z okręgiem i dwa wypełnienia. Jeżeli $\gamma_0,\gamma_1$
są dodatnio zorientowanymi krzywymi brzegowymi
pierścienia, a $h$ jest włóknem, to
$\gamma_0+\gamma_1=0$ w homologii produktu.
Pierwsze wypełnienie narzuca
$\gamma_0+bh=0$, więc daje torus pełny,
którego meridian na pozostałym brzegu ma klasę
$\gamma_1-bh$. Utożsamiamy ten torus pełny
z zewnętrzem węzła trywialnego, przyjmując
bazę chirurgiczną $\mu=h$, $\lambda=\gamma_1-bh$.
Jest ona całkowita, a względem dodatniej bazy
brzegowej $(\gamma_1,h)$ ma wyznacznik $-1$,
zgodnie z orientacją bazy meridian--długość
zewnętrza węzła (normalna zewnętrza jest
przeciwna normalnej jego otoczenia).
Klasa $\lambda$ ogranicza dysk meridianowy
torusa pełnego; identyfikację realizuje
dyfeomorfizm torusów pełnych, jak w dowodzie
lematu \ref{lem:tg-kirby-slam-dunk}.
Drugie wypełnienie ma nachylenie
$\alpha\gamma_1+\beta h$, czyli
\[
(\alpha b+\beta)\mu+\alpha\lambda.
\]
Współczynnik chirurgii wynosi więc
\[
b+\frac{\beta}{\alpha}
=-(p+2)+\frac{p-2}{p-1}
=-\frac{p^2}{p-1}.
\]
To dokładnie $L(p^2,1-p)$ przy konwencji
ustalonej w tezie. Rachunek dotyczy
zorientowanych sklejeń torusów, nie tylko
wyznacznika macierzy przecięcia.
\end{proof}

\begin{przyklad}[Przypadki $p=2,3,4$]
Dla $p=2$ graf składa się z jednego wierzchołka
$-4$, zatem $\partial C_2=L(4,-1)$.
Dla $p=3$ mamy wagi $-5,-2$ i ułamek
$[5,2]^-=5-1/2=9/2$, więc
$\partial C_3=L(9,-2)$.
Dla $p=4$ z wag $-6,-2,-2$ otrzymujemy
$[6,2,2]^-=6-2/3=16/3$ oraz
$\partial C_4=L(16,-3)$.
W każdym przypadku rząd $H_1$ jest kwadratem $p^2$,
ale parametr drugiego argumentu $L(p^2,1-p)$
wymaga pełnej informacji o sklejeniu.
\end{przyklad}

\begin{center}
\begin{tikzpicture}[>=Stealth,scale=.94]
\begin{scope}[shift={(-3.25,0)}]
  \node[font=\small\bfseries,text=tgNavy] at (0,1.10) {$C_2$};
  \node[tgRedVertex] at (0,0) {$-4$};
  \node[font=\scriptsize,text=tgNavy] at (0,-.75) {$L(4,-1)$};
\end{scope}
\begin{scope}[shift={(0,0)}]
  \node[font=\small\bfseries,text=tgNavy] at (0,1.10) {$C_3$};
  \node[tgRedVertex] (c) at (-.8,0) {$-5$};
  \node[tgBlueVertex] (d) at (.8,0) {$-2$};
  \draw[tgTree] (c)--(d);
  \node[font=\scriptsize,text=tgNavy] at (0,-.75) {$L(9,-2)$};
\end{scope}
\begin{scope}[shift={(3.6,0)}]
  \node[font=\small\bfseries,text=tgNavy] at (0,1.10) {$C_4$};
  \node[tgRedVertex] (a) at (-1.25,0) {$-6$};
  \node[tgBlueVertex] (b) at (0,0) {$-2$};
  \node[tgBlueVertex] (c) at (1.25,0) {$-2$};
  \draw[tgTree] (a)--(b)--(c);
  \node[font=\scriptsize,text=tgNavy] at (0,-.75) {$L(16,-3)$};
\end{scope}
\end{tikzpicture}
\end{center}

\section{Wymierna kula homologii Cassona--Harera}
\label{sec:tg-rb-rational-ball}

\begin{definicja}[Wymierna kula homologii]
Zwartą, spójną, zorientowaną czterorozmaitość
$B$ z niepustym brzegiem nazywamy
\emph{wymierną kulą homologii}, jeżeli
$H_*(B;\mathbb Q)\cong H_*(B^4;\mathbb Q)$,
tzn. $H_0(B;\mathbb Q)\cong\mathbb Q$ oraz
$H_k(B;\mathbb Q)=0$ dla każdego $k>0$.
Nie wymagamy, aby $B$ była ściągalna ani
aby $H_1(B;\mathbb Z)$ było zerowe.
\end{definicja}

\begin{twierdzenie}[Casson--Harer; istnienie wypełnienia]
\label{thm:tg-rb-casson-harer}
Dla każdej liczby całkowitej $p\geq2$ istnieje
zwarta, spójna, gładka, zorientowana
czterorozmaitość $B_p$ taka, że (dyfeomorfizm
brzegów zachowuje orientację)
\[
\partial B_p\cong L(p^2,1-p),\qquad
H_1(B_p;\mathbb Z)\cong\mathbb Z/p\mathbb Z,
\quad H_k(B_p;\mathbb Q)=0\ (k>0).
\]
Można ją wybrać z rozkładem na jeden uchwyt
indeksu zero, jeden uchwyt indeksu jeden i jeden
uchwyt indeksu dwa; dla tej konstrukcji
$\pi_1(B_p)\cong\mathbb Z/p\mathbb Z$, a
$\pi_1(\partial B_p)\to\pi_1(B_p)$ jest surjekcją.
\end{twierdzenie}

\begin{proof}[Zakres dowodu i rachunek homologii]
Istnienie takiego \emph{konkretnego rozkładu na uchwyty},
a przede wszystkim identyfikacja jego brzegu jako
$L(p^2,1-p)$, stanowi nietrywialny krok geometryczny.
Korzystamy z lematu 2.1 R.~Fintushela i R.~J. Sterna,
\emph{Rational blowdowns of smooth 4-manifolds},
J. Differential Geom. 46 (1997), 181--235
(arXiv:alg-geom/9505018), gdzie opisano między
innymi konstrukcję Cassona--Harera. Autorzy stosują
inną konwencję oznaczania przestrzeni soczewkowych;
identyfikację orientowanego brzegu w naszej konwencji
podano w twierdzeniu \ref{thm:tg-rb-lens}.
Konkretny diagram uchwytowy $B_p$ przyjmujemy
z rysunku 1 w pracy K.~Yasui,
\emph{Exotic rational elliptic surfaces without
1-handles}, Algebr. Geom. Topol. 8 (2008),
s.~973, sekcja 2 i lemat 2.1. Diagram zawiera
jeden okrąg z kropką (uchwyt indeksu jeden)
i jedną ramowaną składową (uchwyt indeksu dwa)
przechodzącą algebraicznie $p$ razy przez
uchwyt indeksu jeden. Nie wyprowadzamy tu
identyfikacji brzegu tego diagramu z samego
rachunku homologii. Wykażemy natomiast
konsekwencje jego opisu uchwytowego.
Pierwotne źródło istnienia wypełnienia tego typu
to A.~J. Casson, J.~L. Harer,
\emph{Some homology lens spaces which bound
rational homology balls}, Pacific J. Math. 96
(1981), 23--36, twierdzenie na s.~23
i przypadek (b$_1$)(i) na s.~30.
Orientację wypełnienia dobieramy do orientacji
brzegu $C_p$.

W odpowiadającym rozkładowi kompleksie komórkowym
mamy jedną komórkę stopnia zero, jedną stopnia jeden
i jedną stopnia dwa. Krzywa przyczepienia dwukomórki
przebiega algebraicznie $p$ razy przez jedyny
uchwyt indeksu jeden, czyli w homotopii reprezentuje $x^p$
w grupie wolnej na jednym generatorze $x$.
Wartość bezwzględną wykładnika można także
sprawdzić bez odczytywania skrzyżowań diagramu.
Jeżeli wykładnik wynosi $m$, to wymierna
acykliczność wymusza $m\ne0$, a kompleks
komórkowy daje $H_2(B_p;\mathbb Z)=0$,
$H_1(B_p;\mathbb Z)=\mathbb Z/|m|\mathbb Z$
oraz znikanie homologii w stopniach co najmniej trzy.
Z uniwersalnego twierdzenia o współczynnikach
i dualności Poincarégo--Lefschetza otrzymujemy
$H_2(B_p,\partial B_p;\mathbb Z)
\cong\mathbb Z/|m|\mathbb Z$ i
$H_1(B_p,\partial B_p;\mathbb Z)=0$.
Długi ciąg pary daje krótki ciąg dokładny
\[
0\longrightarrow\mathbb Z/|m|\mathbb Z
\longrightarrow H_1(\partial B_p;\mathbb Z)
\longrightarrow\mathbb Z/|m|\mathbb Z
\longrightarrow0.
\]
Ponieważ środkowa grupa ma rząd $p^2$,
z dokładności wynika $m^2=p^2$.
Wybór orientacji komórki pozwala przyjąć $m=p$.
Komórkowy kompleks łańcuchowy w dodatnich stopniach
ma zatem postać
\[
0\longrightarrow C_2=\mathbb Z
\xrightarrow{\ \cdot p\ } C_1=\mathbb Z
\xrightarrow{\ 0\ } C_0=\mathbb Z\longrightarrow0.
\]
Znak pierwszego homomorfizmu ustalamy przez
wybór orientacji komórki. Drugi homomorfizm
jest zerowy, gdyż oba końce jedynej 1-komórki
są dołączone do tej samej 0-komórki.
Mnożenie przez $p$ jest injektywne nad $\mathbb Z$
i odwracalne nad $\mathbb Q$. Wynika stąd $H_2(B_p;\mathbb Z)=0$,
$H_1(B_p;\mathbb Z)\cong\mathbb Z/p\mathbb Z$,
a po rozpatrzeniu tego samego kompleksu łańcuchowego
ze współczynnikami w $\mathbb Q$ również
$H_k(B_p;\mathbb Q)=0$ dla $k>0$.
Ponadto van Kampen daje
$\pi_1(B_p)\cong\langle x\mid x^p=1\rangle$
i stąd grupę cykliczną rzędu $p$.
Aby uzasadnić surjektywność homomorfizmu
$\pi_1(\partial B_p)\to\pi_1(B_p)$,
odwracamy rozkład na uchwyty: względem
kołnierza $\partial B_p$ uchwyty mają
indeksy $4,3,2$, ponieważ uchwyt indeksu
$k$ w rozmaitości czterowymiarowej odpowiada
dualnemu uchwytowi indeksu $4-k$.
Dołączenie uchwytu indeksu dwa narzuca
relację w grupie fundamentalnej (twierdzenie
van Kampena), a uchwyty indeksu co najmniej
trzy jej nie zmieniają: odpowiadają dołączeniu
komórek tych samych wymiarów. Żaden z tych
uchwytów nie wprowadza nowych generatorów.
Ponieważ brzeg jest spójny, otrzymujemy
żądaną surjekcję. Nie wynika ona wyłącznie
z obliczenia pierwszej homologii.
\end{proof}

\begin{uwaga}[Dlaczego $B_p$ nie jest zwykłą kulą]
Dla $p\geq2$ grupa $H_1(B_p;\mathbb Z)$
nie jest zerowa. Zatem $B_p$ ma homologię kuli
wyłącznie ze współczynnikami wymiernymi;
nie jest całkowitą kulą homologii ani
nie może być homeomorficzna z $B^4$.
\end{uwaga}

\section{Forma splecenia i podgrupa metabolizująca}
\label{sec:tg-rb-linkingform}

\begin{stwierdzenie}[Forma splecenia brzegu]
\label{prop:tg-rb-linkingform}
Niech $p\geq2$ i niech $e_1$ będzie generatorem
$H_1(Y_p;\mathbb Z)\cong\mathbb Z/p^2\mathbb Z$
wskazanym w dowodzie twierdzenia
\ref{thm:tg-rb-h1boundary}. Przy konwencji
$\lambda_{Y_p}$ reprezentowanej przez $-Q_p^{-1}$
zachodzi
\[
\lambda_{Y_p}(e_1,e_1)=\frac{p-1}{p^2}
\quad\text{w }\mathbb Q/\mathbb Z.
\]
Podgrupa $p\langle e_1\rangle$ rzędu $p$ jest
izotropowa, a nawet jest równa swemu ortogonalnemu
do niej dopełnieniu.
\end{stwierdzenie}

\begin{proof}
Wzór na element odwrotności macierzy, wyrażony
przez minor, daje
$(Q_p^{-1})_{11}=\det Q_p^{(1)}/\det Q_p$,
gdzie $Q_p^{(1)}$ powstaje przez skreślenie pierwszego
wiersza i kolumny. Dla $p\geq3$ jest to
macierz łańcucha $p-2$ wierzchołków o wadze $-2$,
więc $\det Q_p^{(1)}=(-1)^{p-2}(p-1)$.
Dla $p=2$ używamy konwencji, że wyznacznik
macierzy pustej wynosi jeden. W obu przypadkach
\[
-(Q_p^{-1})_{11}=-\frac{(-1)^{p-2}(p-1)}{(-1)^{p-1}p^2}
=\frac{p-1}{p^2}.
\]
Następnie
$\lambda_{Y_p}(pe_1,pe_1)=p^2(p-1)/p^2=p-1=0$
w $\mathbb Q/\mathbb Z$. Cała podgrupa
$p\langle e_1\rangle$ jest więc izotropowa.
Dla dowolnej klasy $ke_1$ zachodzi
$\lambda_{Y_p}(ke_1,pe_1)=k(p-1)/p$
w $\mathbb Q/\mathbb Z$.
Wartość ta jest zerowa wtedy i tylko wtedy,
gdy $p\mid k(p-1)$, a ponieważ
$\gcd(p,p-1)=1$, jest to równoważne $p\mid k$.
Stąd dopełnienie ortogonalne jest dokładnie
$p\langle e_1\rangle$. W szczególności
równość zachodzi również dla złożonych $p$.
\end{proof}

\begin{uwaga}[Związek z wypełnieniem]
Surjekcja $H_1(Y_p;\mathbb Z)\to H_1(B_p;\mathbb Z)$
ma jądro rzędu $p$; jest to właśnie
$p\langle e_1\rangle$.
Taki układ jest zgodny z ogólną zasadą,
że jądro mapy brzegowej dla wymiernej kuli homologii
jest metabolizerem formy splecenia.
W szczególności istnienie $B_p$ nakłada ograniczenie
na brzeg znacznie silniejsze niż samo $|H_1(Y_p)|=p^2$.
\end{uwaga}

\begin{stwierdzenie}[Dokładny ciąg homologii wymiernej kuli]
\label{st:tg-rb-boundary-kernel}
Niech $p\geq2$, a $B_p$ niech będzie
standardową wymierną kulą homologii
z Twierdzenia~\ref{thm:tg-rb-casson-harer},
otrzymaną z uchwytów indeksów $0,1,2$.
Niech $Y_p=\partial B_p$.
Wówczas istnieje dokładny ciąg
\[
0\longrightarrow\mathbb Z/p\mathbb Z
\longrightarrow H_1(Y_p;\mathbb Z)
\longrightarrow H_1(B_p;\mathbb Z)
\longrightarrow0.
\]
Dla identyfikacji $H_1(Y_p;\mathbb Z)
\cong\mathbb Z/p^2\mathbb Z$ jego pierwszy wyraz
jest jedyną podgrupą rzędu $p$,
a więc $pH_1(Y_p;\mathbb Z)$.
Jest to metabolizer formy splecenia
$\lambda_{Y_p}$.
\end{stwierdzenie}

\begin{proof}
Z opisu uchwytowego $B_p$ wynika
$H_1(B_p;\mathbb Z)=\mathbb Z/p\mathbb Z$
oraz $H_2(B_p;\mathbb Z)=H_3(B_p;\mathbb Z)=0$.
Z uniwersalnego twierdzenia o współczynnikach:
\[
H^2(B_p;\mathbb Z)
\cong\operatorname{Ext}^1_{\mathbb Z}
(H_1(B_p;\mathbb Z),\mathbb Z)
\cong\mathbb Z/p\mathbb Z,
\qquad
H^3(B_p;\mathbb Z)=0.
\]
Dualność Poincarégo--Lefschetza daje zatem
$H_2(B_p,Y_p;\mathbb Z)\cong\mathbb Z/p\mathbb Z$
i $H_1(B_p,Y_p;\mathbb Z)=0$.
Odpowiedni fragment długiego ciągu pary
ma postać
\[
H_2(B_p;\mathbb Z)\longrightarrow
H_2(B_p,Y_p;\mathbb Z)\longrightarrow
H_1(Y_p;\mathbb Z)\longrightarrow
H_1(B_p;\mathbb Z)\longrightarrow
H_1(B_p,Y_p;\mathbb Z).
\]
Pierwsza i ostatnia grupa znikają,
co daje dokładny ciąg z tezy.
Ponieważ grupa środkowa jest cykliczna
rzędu $p^2$, jej jedyna podgrupa
rzędu $p$ to $pH_1(Y_p;\mathbb Z)$.
Poprzednie stwierdzenie
\ref{prop:tg-rb-linkingform}
wykazało bezpośrednio, że jest
ona izotropowa i równa swemu
dopełnieniu ortogonalnemu.
Wobec tego stanowi metabolizer.
\end{proof}

\section{Definicja operacji i zmiana niezmienników}
\label{sec:tg-rb-definition}

\begin{definicja}[Wymierna kontrakcja]
\label{def:tg-rb-operation}
Niech $p\geq2$, a $X$ niech będzie zamkniętą,
spójną, zorientowaną, gładką czterorozmaitością.
Załóżmy, że $X$ zawiera w swoim wnętrzu
gładko zanurzoną kopię $C_p$ jako podrozmaitość
z brzegiem, przez zanurzenie zachowujące orientację.
Ustalamy tę identyfikację $C_p$ z jego obrazem. Oznaczmy
$W:=X\setminus\operatorname{int}C_p$ i
$Y_p:=\partial C_p$.
\emph{Wymierną kontrakcją $X$ wzdłuż $C_p$}
nazywamy rozmaitość
\[
X_{(p)}:=W\cup_{\varphi}B_p,
\]
gdzie $\varphi:\partial B_p\to\partial C_p$ jest
dyfeomorfizmem \emph{zachowującym} orientację
brzegową $\partial C_p$ i $\partial B_p$.
Ponieważ $\partial W=-\partial C_p$, to po
utożsamieniu $\partial W$ z $\partial C_p$
to samo sklejenie jest odwracające orientację
między brzegami $W$ i $B_p$.
W klasycznym modelu $B_p$ typ otrzymanej
gładkiej rozmaitości jest niezależny od wyboru
takiej identyfikacji; korzystamy z twierdzenia
o przedłużaniu zachowujących orientację
samodyfeomorfizmów $\partial B_p$ do $B_p$
(R.~Fintushel, R.~J. Stern,
\emph{Rational blowdowns of smooth 4-manifolds},
J. Differential Geom. 46 (1997), 181--235, wniosek 2.2).
Istotnie, dla dwóch sklejeń $\varphi,\psi$
samodyfeomorfizm $\psi^{-1}\varphi$ brzegu $B_p$
przedłuża się do dyfeomorfizmu $F$ kuli $B_p$.
Po uzgodnieniu kołnierzy odwzorowanie równe
identyczności na $W$ i $F$ na $B_p$ daje
dyfeomorfizm $W\cup_\varphi B_p\to W\cup_\psi B_p$.
W dalszych rachunkach wystarcza jednak ustalenie
jednego sklejenia, bez stosowania tego twierdzenia.
\end{definicja}

\begin{center}
\begin{tikzpicture}[>=Stealth,scale=.92]
  \node[draw=tgNavy,fill=tgNavy!6,rounded corners=3pt,
    minimum width=23mm,minimum height=13mm,align=center,
    font=\footnotesize] (W) at (0,1.7) {$W$\\{\scriptsize wspólne dopełnienie}};
  \node[draw=tgCoral,fill=tgCoral!14,rounded corners=3pt,
    minimum width=27mm,minimum height=13mm,align=center,
    font=\footnotesize] (C) at (-3.0,0) {$C_p$\\{\scriptsize plumbing sfer}};
  \node[draw=tgTeal,fill=tgTeal!14,rounded corners=3pt,
    minimum width=27mm,minimum height=13mm,align=center,
    font=\footnotesize] (B) at (3.0,0) {$B_p$\\{\scriptsize wymierna kula}};
  \draw[tgAccent,-{Stealth[length=3mm]}] (W.south west)--(C.north east);
  \draw[tgTree,-{Stealth[length=3mm]}] (W.south east)--(B.north west);
  \node[font=\footnotesize,text=tgCoral] at (-3.0,-1.05) {$X=W\cup_{Y_p}C_p$};
  \node[font=\footnotesize,text=tgTeal] at (3.0,-1.05) {$X_{(p)}=W\cup_{Y_p}B_p$};
  \node[font=\scriptsize,text=tgNavy] at (0,-1.65)
    {Ten sam brzeg $Y_p\cong L(p^2,1-p)$, różne wypełnienia.};
\end{tikzpicture}
\end{center}

\begin{twierdzenie}[Skutek liczbowy wymiernej kontrakcji]
\label{thm:tg-rb-invariants}
Niech $p\geq2$ i niech $X$, $X_{(p)}$ oraz
$C_p\subset X$ będą jak w definicji
\ref{def:tg-rb-operation}.
Niech $e$ oznacza charakterystykę Eulera,
$\sigma$ sygnaturę formy przecięcia, a
$b_2^\pm$ wymiary maksymalnych dodatnio i
ujemnie określonych podprzestrzeni
$H_2(-;\mathbb R)$.
Wtedy
\[
\begin{aligned}
e(X_{(p)})&=e(X)-(p-1), &
\sigma(X_{(p)})&=\sigma(X)+(p-1),\\
b_2^+(X_{(p)})&=b_2^+(X), &
b_2^-(X_{(p)})&=b_2^-(X)-(p-1).
\end{aligned}
\]
W szczególności, jeśli wprowadzimy oznaczenia
$\chi_h(Z):=(e(Z)+\sigma(Z))/4$ oraz
$c_{\mathrm{num}}(Z):=2e(Z)+3\sigma(Z)$,
to $\chi_h(X_{(p)})=\chi_h(X)$ i
$c_{\mathrm{num}}(X_{(p)})=c_{\mathrm{num}}(X)+(p-1)$.
Dla zamkniętych czterorozmaitości prawie zespolonych,
orientowanych przez strukturę prawie zespoloną,
ostatnia wielkość jest liczbą Chernowską $c_1^2$.
Nie zakładamy tu, że każda rozmaitość występująca
w twierdzeniu jest prawie zespolona. W takim przypadku
$\chi_h$ jest tylko oznaczeniem liczbowym, a nie holomorficzną
charakterystyką Eulera.
\end{twierdzenie}

\begin{proof}
Brzeg $Y_p$ jest wymierną sferą homologiczną,
więc $H_1(Y_p;\mathbb Q)=H_2(Y_p;\mathbb Q)=0$.
Rozpatrzmy rozkłady $X=W\cup C_p$ oraz
$X_{(p)}=W\cup B_p$. Ciągi Mayera--Vietorisa
nad $\mathbb Q$ dają izomorfizmy
\[
H_2(X;\mathbb Q)\cong H_2(W;\mathbb Q)
\oplus H_2(C_p;\mathbb Q),\qquad
H_2(X_{(p)};\mathbb Q)\cong H_2(W;\mathbb Q).
\]
Izomorfizmy te pochodzą od włożeń składników:
w ciągu Mayera--Vietorisa zarówno grupa
$H_2(Y_p;\mathbb Q)$ poprzedzająca sumę,
jak i grupa $H_1(Y_p;\mathbb Q)$ następująca
po niej znikają. Dzięki kołnierzom brzegowym
cykle z poszczególnych składników można przesunąć
do ich wnętrz i obliczać przecięcia w tych wnętrzach.
Reprezentanci z różnych składników są rozłączni,
toteż pierwsza suma jest ortogonalna względem
formy przecięcia. Forma na wspólnym składniku $W$
jest w obu rozkładach ta sama i jest nieosobliwa:
w ciągu pary $(W,\partial W)$ grupy
$H_2(\partial W;\mathbb Q)$ i
$H_1(\partial W;\mathbb Q)$ znikają, więc
$H_2(W;\mathbb Q)\to H_2(W,\partial W;\mathbb Q)$
jest izomorfizmem. Dualność Poincarégo--Lefschetza
identyfikuje grupę względną z przestrzenią
sprzężoną do $H_2(W;\mathbb Q)$, dając
nieosobliwość formy. Rozszerzenie skalarów
do $\mathbb R$ zachowuje te rozkłady.
Na $H_2(C_p;\mathbb Q)$ forma jest ujemnie określona
rangi $p-1$, jak wykazaliśmy w lemacie
\ref{lem:tg-rb-minors}. Otrzymujemy stąd równości
dla $b_2^\pm$ oraz dla sygnatury.

Wzór na charakterystykę Eulera korzysta z addytywności
$e(A\cup_YB)=e(A)+e(B)-e(Y)$ dla zwartych
kompleksów triangulowalnych. Ponieważ
$e(Y_p)=0$, $e(C_p)=1+(p-1)=p$ oraz
$e(B_p)=1$, mamy
$e(X_{(p)})-e(X)=e(B_p)-e(C_p)=1-p$.
Równania dla $\chi_h$ i $c_{\mathrm{num}}$ są
następnie bezpośrednim podstawieniem:
$\Delta(e+\sigma)=0$ oraz
$\Delta(2e+3\sigma)=2(1-p)+3(p-1)=p-1$.
Interpretacja ostatniej wielkości w przypadku
struktury prawie zespolonej wynika z
$c_2[Z]=e(Z)$ i $p_1=c_1^2-2c_2$
dla rzeczywistej wiązki stycznej uzyskanej
z zespolonej wiązki rzędu dwa, a następnie
z twierdzenia Hirzebrucha o sygnaturze
$p_1[Z]=3\sigma(Z)$.
\end{proof}

\begin{stwierdzenie}[Pierwsza homologia całkowita po wymiernej kontrakcji]
\label{prop:tg-rb-integral-h1}
Przy założeniach definicji \ref{def:tg-rb-operation}
niech $a\in H_1(W;\mathbb Z)$ będzie obrazem generatora
$H_1(Y_p;\mathbb Z)\cong\mathbb Z/p^2\mathbb Z$.
Wówczas
\[
H_1(X;\mathbb Z)\cong H_1(W;\mathbb Z)/\langle a\rangle,
\qquad
H_1(X_{(p)};\mathbb Z)\cong H_1(W;\mathbb Z)/\langle pa\rangle.
\]
Naturalne odwzorowanie drugiego ilorazu na pierwszy
jest surjekcją, której jądro jest cykliczne
o rzędzie dzielącym $p$.
\end{stwierdzenie}
\begin{proof}
Brzeg $Y_p$ jest spójny. Każda składowa $W$
spotyka $Y_p$: gdyby nie spotykała brzegu,
byłaby zarazem otwarta i domknięta w spójnej
rozmaitości $X$ oraz rozłączna z $C_p$.
Ponieważ spójny brzeg leży w jednej składowej
$W$, dopełnienie $W$ jest spójne. Obie
dołączane części są również spójne,
więc w obu ciągach Mayera--Vietorisa mapa
w stopniu zero z homologii wspólnego brzegu
do sumy homologii składników jest injektywna.
Ponieważ $H_1(C_p;\mathbb Z)=0$, w pierwszym
rozkładzie otrzymujemy iloraz przez $\langle a\rangle$.
W drugim rozkładzie mamy koker mapy
$H_1(Y_p;\mathbb Z)\to H_1(W;\mathbb Z)
\oplus H_1(B_p;\mathbb Z)$.
Ze stwierdzenia \ref{st:tg-rb-boundary-kernel}
wynika, że obraz generatora w drugim składniku jest
generatorem $b$ grupy cyklicznej rzędu $p$.
Zatem koker ma prezentację
\[
\bigl(H_1(W;\mathbb Z)\oplus\langle b\mid pb=0\rangle\bigr)
 /\langle(a,-b)\rangle .
\]
Relacja utożsamia $b$ z $a$, a $pb=0$
przechodzi w $pa=0$. Jawnie odwzorowanie
$(x,kb)\mapsto[x+ka]$ do ilorazu przez $\langle pa\rangle$
jest dobrze określone i znika na relacji $(a,-b)$.
Jego odwrotność jest indukowana przez
$x\mapsto(x,0)$, ponieważ klasy $(0,b)$
i $(a,0)$ są równe w kokerze.
Jądro naturalnej surjekcji ilorazów to
$\langle a\rangle/\langle pa\rangle$,
cykliczne i anihilowane przez $p$.
W szczególności liczby Bettiego w stopniu jeden
są równe, choć grupy całkowite mogą się różnić.
\end{proof}

\begin{przyklad}[Kontrola dla $p=2$ i $p=3$]
Dla $p=2$ usuwamy jedną sferę o samoprzecięciu
$-4$. Ponieważ $e(C_2)=2$, $\sigma(C_2)=-1$,
a $e(B_2)=1$, $\sigma(B_2)=0$, otrzymujemy
$\Delta e=-1$ i $\Delta\sigma=1$.
Dla $p=3$ usuwamy łańcuch sfer o wagach $-5,-2$:
$e(C_3)=3$, $\sigma(C_3)=-2$, zatem
$\Delta e=-2$ i $\Delta\sigma=2$.
W obu przypadkach wzrost $2e+3\sigma$ wynosi
odpowiednio $1$ i $2$.
\end{przyklad}

\section{Grupa fundamentalna po wymiernej kontrakcji}
\label{sec:tg-rb-pi1}

\begin{stwierdzenie}[Kryterium prostej spójności]
\label{prop:tg-rb-pi1}
Dla $p\geq2$, przy założeniach definicji
\ref{def:tg-rb-operation},
jeżeli $W:=X\setminus\operatorname{int}C_p$ jest
po prostu spójna, to $X_{(p)}$ jest po prostu spójna.
\end{stwierdzenie}

\begin{proof}
Wybierzmy punkt bazowy na wspólnym brzegu $Y_p$.
Brzeg $Y_p$ jest spójny. Także $W$ jest spójna:
każda składowa $W$ musi spotykać $\partial W$,
gdyż w przeciwnym razie stanowiłaby otwartą
i domkniętą składową $X$ rozłączną z $C_p$.
Spójny brzeg należy do jednej składowej $W$.
Rozszerzamy oba składniki przez otwarte kołnierze,
tak aby ich przecięcie retraktowało się na $Y_p$,
a same rozszerzenia na odpowiednio $W$ i $B_p$.
Twierdzenie van Kampena daje prezentację
\[
\pi_1(X_{(p)})\cong\pi_1(W)*_{\pi_1(Y_p)}\pi_1(B_p).
\]
Jeśli $\pi_1(W)=1$, pozostaje iloraz
$\pi_1(B_p)$ przez normalne domknięcie obrazu
$\pi_1(Y_p)\to\pi_1(B_p)$. Na mocy twierdzenia
Cassona--Harera obraz ten jest całą grupą
$\pi_1(B_p)\cong\mathbb Z/p\mathbb Z$,
więc iloraz jest trywialny.
\end{proof}

\begin{stwierdzenie}[Ogólny rachunek grupy po sklejeniu]
\label{prop:tg-rb-pi1-general}
Niech $p\ge2$, $C_p\subset X$ i $B_p$
będą jak w definicji wymiernej kontrakcji,
a $W=X\setminus\operatorname{int}C_p$.
Oznaczmy przez $Y_p$ wspólny brzeg oraz
przez $i_W\colon\pi_1(Y_p)\to\pi_1(W)$
i $i_B\colon\pi_1(Y_p)\to\pi_1(B_p)$
homomorfizmy indukowane przez włożenia,
przy identyfikacji brzegów i zgodnym
wyborze punktu bazowego.
Ponieważ $i_B$ jest surjektywny, zachodzi
\[
\pi_1(X_{(p)})\cong
\pi_1(W)\big/
\big\langle\!\big\langle
i_W(\ker i_B)
\big\rangle\!\big\rangle,
\]
gdzie podwójne nawiasy oznaczają
domknięcie normalne podzbioru
w grupie $\pi_1(W)$.
W szczególności $X_{(p)}$ jest
po prostu spójna wtedy i tylko wtedy, gdy
$i_W(\ker i_B)$ normalnie generuje
całą grupę $\pi_1(W)$.
\end{stwierdzenie}

\begin{proof}
Z twierdzenia van Kampena
$\pi_1(X_{(p)})$ jest ilorazem
iloczynu wolnego $\pi_1(W)*\pi_1(B_p)$
przez relacje $i_W(\gamma)=i_B(\gamma)$
dla wszystkich $\gamma\in\pi_1(Y_p)$.
Oznaczmy ten iloraz iloczynu wolnego przez $G$,
a przez $N$ normalne domknięcie
$i_W(\ker i_B)$ w $\pi_1(W)$.
Naturalny homomorfizm $\pi_1(W)\to G$
jest surjektywny, ponieważ każdy element
$\pi_1(B_p)$ ma postać $i_B(\gamma)$
i w $G$ jest utożsamiany z $i_W(\gamma)$.
Ponadto zabija $N$, więc indukuje
surjekcję $\pi_1(W)/N\to G$.
Aby skonstruować jej odwrotność, oznaczmy
przez $q$ projekcję $\pi_1(W)\to\pi_1(W)/N$.
Dla $b\in\pi_1(B_p)$ wybieramy $\gamma$
z $i_B(\gamma)=b$ i kładziemy
$f(b)=q(i_W(\gamma))$.
Jeżeli wybierzemy inny przeciwobraz $\gamma'$,
to $\gamma'^{-1}\gamma\in\ker i_B$, więc
$q(i_W(\gamma'))=q(i_W(\gamma))$.
Odwzorowanie $f$ jest dobrze określonym
homomorfizmem, co wynika z wyboru
iloczynu przeciwobrazów dla iloczynu w grupie.
Zachodzi $f\circ i_B=q\circ i_W$.
Własność uniwersalna iloczynu wolnego
z amalgamacją daje homomorfizm
$G\to\pi_1(W)/N$. Oba złożenia są
identycznościami na generatorach odpowiednich grup,
co dowodzi wskazanego izomorfizmu.
Jest on trywialny dokładnie wtedy,
gdy domknięcie normalne jest
całą grupą $\pi_1(W)$.
Wcześniejsze kryterium stanowi
przypadek szczególny dla
$\pi_1(W)=1$.
\end{proof}

\begin{uwaga}[Dlaczego $\pi_1(X)=1$ nie wystarcza w dowodzie]
Z prostej spójności $X$ nie wynika automatycznie,
że $W=X\setminus\operatorname{int}C_p$ jest po prostu
spójna. Ponieważ $C_p$ jest po prostu spójna
(ma rozkład z uchwytami indeksów zero i dwa),
analogiczny rachunek van Kampena daje
$\pi_1(X)\cong\pi_1(W)/
\langle\!\langle i_W(\pi_1(Y_p))\rangle\!\rangle$.
Jeżeli $\pi_1(X)=1$, to obraz całej grupy brzegowej
normalnie generuje $\pi_1(W)$. Nie oznacza to,
że już obraz jej właściwej podgrupy
$\ker i_B$ normalnie generuje tę grupę.
Właśnie tego wymaga stwierdzenie
\ref{prop:tg-rb-pi1-general}.
Grupy przy sklejaniu należy zatem obliczyć oddzielnie.
W konstrukcjach egzotycznych czterorozmaitości
jest to jeden z istotnych etapów dowodu, którego
nie wolno zastąpić rachunkiem samej homologii.
\end{uwaga}

\section{Zastosowanie do egzotycznych struktur gładkich}
\label{sec:tg-rb-exotic}

\begin{definicja}[Egzotyczna struktura gładka]
Dwie zamknięte, spójne, zorientowane, gładkie
czterorozmaitości $X_1,X_2$ nazywamy
\emph{egzotyczną parą}, jeżeli są homeomorficzne
z zachowaniem orientacji, ale nie są dyfeomorficzne.
\end{definicja}

Sama wymierna kontrakcja nie dowodzi egzotyczności.
Trzeba wykonać trzy niezależne czynności:
obliczyć $\pi_1$ nowej rozmaitości, ustalić
jej formę przecięcia oraz znaleźć niezmiennik
\emph{gładki}, który odróżnia wynik od znanego modelu.
Twierdzenie Freedmana służy do klasyfikacji
topologicznej po spełnieniu właściwych hipotez;
niezmienniki Seiberga--Wittena lub Donaldsona
służą do rozróżniania struktur gładkich.

\begin{przyklad}[Dane modelu $\mathbb{CP}^2\#7\overline{\mathbb{CP}}^{\,2}$]
Niech $Z=\mathbb{CP}^2\#7\overline{\mathbb{CP}}^{\,2}$
z orientacją standardową na pierwszym składniku.
Z rozkładu sumy spójnej mamy
$H_2(Z;\mathbb Z)\cong\mathbb Z^8$ oraz
\[
Q_Z=\langle1\rangle\oplus7\langle-1\rangle,
\qquad b_2^+(Z)=1,\quad b_2^-(Z)=7.
\]
Ponieważ $Z$ jest po prostu spójna,
$e(Z)=2+b_2(Z)=10$, a
$\sigma(Z)=1-7=-6$. W szczególności
$2e(Z)+3\sigma(Z)=20-18=2$ i
$\chi_h(Z)=(10-6)/4=1$.

Te liczby są \emph{konieczne}, lecz nie wystarczają,
aby dowolna rozmaitość była homeomorficzna z $Z$.
Wymagana jest m.in. prosta spójność i ustalenie
całkowitej formy przecięcia wraz z jej parzystością.
Przy tych założeniach twierdzenie Freedmana
może dać klasyfikację topologiczną, lecz nie
rozstrzyga problemu dyfeomorfizmu.
\end{przyklad}

\begin{twierdzenie}[Park; zastosowanie wymiernej kontrakcji,
przyjmowane z literatury]
Istnieje zamknięta, po prostu spójna, gładka
czterorozmaitość symplektyczna $X_{\mathrm{Park}}$
o $b_2^+=1$ i $c_1^2=2$,
homeomorficzna z
$\mathbb{CP}^2\#7\overline{\mathbb{CP}}^{\,2}$,
lecz z nią niedyfeomorficzna. W konstrukcji
wykorzystuje się m.in. wymierną kontrakcję
odpowiedniego plumbingowania oraz niezmienniki
Seiberga--Wittena.
\end{twierdzenie}

\begin{uwaga}[Co w tym przykładzie pozostaje bez dowodu]
Powyższe twierdzenie jest wynikiem J.~Parka,
\emph{Simply connected symplectic 4-manifolds
with $b_2^+=1$ and $c_1^2=2$},
Invent. Math. 159 (2005), 657--667,
twierdzenie 1.1 (arXiv:math/0311395).
Nasz rachunek $e$, $\sigma$ i $Q$ \emph{nie} jest
jego dowodem. Pełna konstrukcja wymaga znalezienia
odpowiedniego zanurzenia konfiguracji w wyjściowej
rozmaitości, sprawdzenia grupy fundamentalnej
dopełnienia oraz obliczenia gładkich niezmienników.
Ponadto dla $b_2^+=1$ niezmienniki
Seiberga--Wittena zależą od wyboru komory;
porównanie z rozmaitościami wymiernymi wymaga
konsekwentnego ustalenia tej komory.
\end{uwaga}

\begin{twierdzenie}[Symington; wersja symplektyczna, bez dowodu]
Niech $(X,\omega)$ będzie zamkniętą symplektyczną
czterorozmaitością zawierającą symplektycznie
zanurzony standardowy model plumbingowania
$C_p$, w którym sfery są symplektyczne, a sąsiednie
sfery przecinają się ortogonalnie względem formy
symplektycznej. Wtedy istnieje struktura
symplektyczna na $B_p$ i zgodne sklejenie
symplektyczne w kołnierzach, dające formę
symplektyczną na rozmaitości po wymiernej
kontrakcji, zgodną z pierwotną formą na
wspólnym dopełnieniu.
\end{twierdzenie}

Przyjmujemy ten wynik z pracy M.~Symington,
\emph{Symplectic rational blowdowns},
J. Differential Geom. 50 (1998), 505--518,
twierdzenie 1.3 i uwaga 1.2
(arXiv:math/9802079). Standardowy model można
uzyskać z konfiguracji symplektycznych sfer
o odpowiednim wzorze przecięć, po symplektycznej
izotopii sfer zapewniającej ortogonalność
przecięć i zastosowaniu twierdzenia o otoczeniu
symplektycznym. Istnienie zgodnego sklejenia
stanowi część tezy twierdzenia, nie wynika
z samej dyfeomorfii orientowanych brzegów.
Nie twierdzimy, że dowolnie wybrane sklejenie
zachowuje z góry zadane formy symplektyczne.
W ogólniejszym sformułowaniu wystarcza
zanurzenie gładkiego plumbingu, którego sfery
są symplektyczne, z dodatnimi transwersalnymi
przecięciami sąsiednich sfer. Zob.
M.~Symington, \emph{Generalized symplectic
rational blowdowns}, Algebr. Geom. Topol. 1
(2001), 503--518, twierdzenie 1.3 na s.~505
(arXiv:math/0108075), dla parametru $q=1$.
W sekcji 4 autorka uzyskuje zgodność
kołnierzy z lokalnych modeli symplektycznych.
Nie wynika stąd jednoznaczność uzyskanej
formy z dokładnością do symplektomorfizmu.

\section{Przejście do uogólnień}
\label{sec:tg-rb-general}
Gdy $p>q\geq1$ są względnie pierwsze, analogiczne
plumbingowanie $C_{p,q}$ wyznacza ujemne rozwinięcie
$\frac{p^2}{pq-1}$. Rachunek grupy brzegowej
$H_1(\partial C_{p,q};\mathbb Z)\cong\mathbb Z/p^2\mathbb Z$
nie wystarcza do identyfikacji jej orientowanego typu.
W następnym rozdziale wyprowadzimy ujemne ułamki łańcuchowe,
pełny rachunek formy przecięcia i rzeczywisty współczynnik
chirurgii, a następnie zastosujemy je do wymiernych kontrakcji.
Aby sprawdzić, kiedy otrzymujemy nową strukturę gładką,
potrzebujemy jeszcze niezmienników Seiberga--Wittena.

% Iteracja 10: uogolniona wymierna kontrakcja i podstawy niezmiennikow SW.
\begin{uwaga}[Granica rachunku wyznacznika]
Sama tożsamość $|\det Q_{C_p}|=p^2$ wyznacza rząd
grupy $H_1(\partial C_p;\mathbb Z)$, ale nie rozstrzyga
ani izomorfizmu tej grupy z $\mathbb Z/p^2\mathbb Z$,
ani typu dyfeomorfizmu brzegu. Cykliczność trzeba
wyprowadzić z prezentacji relacji lub jej postaci Smitha,
a identyfikację przestrzeni soczewkowej z odpowiedniego
ułamka łańcuchowego. Istnienie zastępującej ją wymiernej
kuli homologii jest kolejnym, niezależnym wynikiem
geometrycznym (Casson--Harer).
\end{uwaga}

